This Delivery Note accompanies \texttt{\deliveryid}, a set of conflict forecasts produced by eight
HydraNet models on the calibration partition at global PRIO-GRID month, each collapsed from its
posterior draws three different ways. It is the authoritative description of what was handed over.
The folder was uploaded on 29 September 2026 with a short README and an audit notebook and no note;
this document supplies the note, and supersedes that README where the two disagree.

The purpose of this note is to:
\begin{itemize}
    \item record what was delivered, in enough detail to identify it again later;
    \item describe the models that produced it, and what distinguishes them from one another;
    \item explain why there are three files per model per origin rather than one, and what each
          one is;
    \item state how the files are organised and how they are intended to be consumed; and
    \item set out what is \emph{not} in the delivery, and what is not yet known about it.
\end{itemize}

\noindent This note covers:
\begin{itemize}
    \item the architecture all eight share, and the eight configurations that differ
          (Sections~\ref{sec:model}--\ref{sec:models});
    \item the three point estimators, the arithmetic of each, and which one the models declare
          (Section~\ref{sec:frames});
    \item the file layout, schema and units of the delivered forecasts
          (Section~\ref{sec:files});
    \item measured magnitude and horizon behaviour --- diagnostics, not an evaluation
          (Section~\ref{sec:magnitude}); and
    \item what the delivery excludes, and the caveats that bear on using it
          (Sections~\ref{sec:excluded}--\ref{sec:caveats}).
\end{itemize}

\noindent This note explicitly does \emph{not} cover:
\begin{itemize}
    \item \textbf{any evaluation against observed outcomes.} Section~\ref{sec:magnitude} reports
          what the models predict and how that varies by forecast step and by estimator. It scores
          nothing, ranks nothing and compares nothing to actuals beyond the single magnitude ratio
          the run itself computed;
    \item any recommendation as to which of the eight, or which of the three estimators, to
          prefer. That is what the selection routine is for;
    \item the validation or forecasting partitions, or any live forecast; and
    \item the internal machinery that produced the run --- the cost, the hardware and what the
          deployment cost us are recorded separately, in
          \texttt{reports/runpod\_cost\_and\_time\_note\_2026-09.md} and
          \texttt{reports/postmortem\_runpod\_first\_deployment\_2026-09.md}.
\end{itemize}

\noindent The delivery is a fixed artefact. It is identified by the identifier and checksum on
the title page --- quote both --- and is not revised in place: any correction is issued as a new
delivery with a new identifier and a new checksum.
